The Clifford-module restriction gives an exact numerical range for
\eqref{eq:orthogonal-quartic-data}. We use the conventional function
\[
 (\delta(1),\ldots,\delta(8))=(1,2,4,4,8,8,8,8),
 \qquad \delta(k+8)=16\delta(k).
\]
Here $\delta(k)$ is the real irreducible module dimension for the
negative Clifford algebra on $k-1$ generators. A symmetric Clifford
system with $p\ge2$ matrices on $\R^{2\ell}$ exists exactly when
$\delta(p-1)\mid\ell$ \cite[Section~3.5]{FKM81}.
See also \cite[Chapter~0, pp.~3--4]{Shapiro2000} for the equivalent
Hurwitz--Radon restriction.

\begin{proposition}[Realizable quartic weights]
\label{prop:quartic-weight-realizability}
Let $r_0,r_1,r_2$ be positive integers, write their rearrangement as
$a\le b\le c$, and set $\ell=a+b+c-1$.
Nonzero quartics satisfying \eqref{eq:orthogonal-quartic-data} on
$\R^{2\ell}$ exist if and only if \(\delta(a+b-1)\mid\ell\).
Whenever this condition holds, a realization is obtained from one
Clifford system by two blocks of squared quadratic coordinates and
their complement in $R^2$.
\end{proposition}

\begin{proof}[Construction and parameter restrictions]
\leavevmode\par\nobreak
Assume the divisibility condition. Take $a+b$ symmetric Clifford
matrices, with quadratic coordinates $p_1,\ldots,p_{a+b}$.
The three polynomials
\[
 \sum_{\alpha=1}^{a}p_\alpha^2,\qquad
 \sum_{\alpha=a+1}^{a+b}p_\alpha^2,\qquad
 R^2-\sum_{\alpha=1}^{a+b}p_\alpha^2
\]
have weights $(a,b,c)$ by \cref{lem:Clifford-block-quartics}.
Their positive Laplacian coefficients ensure that none is zero.

For any realization, summing the Laplacians forces
$N=2(r_0+r_1+r_2-1)$. Positivity and integrality of the weights follow from
\cref{lem:quartic-block-weights-image}.
The converse uses the completed classification with four principal
curvatures and compatibility of the three blocks. See
\cref{app:quartic-weight-realizability} and its weight-one,
Clifford-joining, and exceptional-triple lemmas
\cref{lem:quartic-weight-one,lem:quartic-Clifford-joining,lem:quartic-exclude-155}.
\end{proof}

This is a classification of realizable weights, not of simultaneous
orthogonal equivalence classes of triples. For example, $(3,3,5)$
is excluded since $\delta(5)=8\nmid10$, although the individual
Cartan--M\"unzner multiplicity pairs $(2,7),(2,7),(4,5)$ all occur.
The criterion is symmetric in the weights; the subcalibration
range in \cref{thm:quartic-table-existence} instead depends on
which two blocks are chosen for the defining polynomial.

\begin{corollary}[The weights with active sum nine]
\label{cor:quartic-nine-weights}
Let $r_0,r_1,r_2$ be positive integers with $r_1,r_2\geq2$ and
$r_1+r_2=9$. A quartic realization exists if and only if $r_0$ is
a positive multiple of eight. Its ambient dimension is $N=2r_0+16$.
\end{corollary}

\begin{proof}
Up to exchanging the active weights, their pairs are $(2,7)$,
$(3,6)$, and $(4,5)$. For $1\leq r_0\leq7$, none of the corresponding module
dimensions divides $\ell=r_0+8$; the seven checks are listed in
\cref{tab:quartic-nine-low-weights}.
If $r_0\geq8$, the two
smallest weights are the active pair. Hence $a+b-1=8$ and the same
criterion is $8=\delta(8)\mid r_0+8$, equivalently $8\mid r_0$.
\end{proof}
